\documentclass[11pt]{article}
\usepackage[margin=1in]{geometry}
\usepackage{amsmath,amssymb,amsthm}
\usepackage{graphicx}
\usepackage{booktabs}
\usepackage{listings}
\usepackage{natbib}
\usepackage[hidelinks]{hyperref}
\newtheorem{theorem}{Theorem}
\newtheorem{lemma}[theorem]{Lemma}

\title{ARGO: A Provenance-Preserving Research Architecture\\
for Autonomous Empirical Agents}
\author{Yunsang Um}
\date{Evidence cutoff: 2 September 2026}

\begin{document}
\maketitle

\begin{abstract}
Autonomous research systems can propose hypotheses, edit code, run experiments,
and draft papers, but these abilities do not by themselves establish scientific
validity.  A long-running agent may silently change a comparison after observing
an outcome, detach a claim from the artifact that supports it, confuse execution
failure with scientific evidence, or present an inherited mechanism as a new
contribution.  This thesis specifies candidate requirements for ARGO, a
research-agent architecture intended to make those failures observable and
rejectable.  ARGO is designed as a progressive migration of a Prime Agent fork;
it is not a wrapper around Prime Agent and not a clean-room origin story.  It
inherits Prime Agent's persistent computation, daemon-backed continuity,
recursive subagents, messaging, recovery, and Continual Harness.  Proposed ARGO
deltas concern scientific decision records, three-part comparison identity,
claim-typed evidence admission, separation of research refinement from engine
refinement, and evidence-bound manuscript builds.

I evaluate the design prospectively rather than reconstructing a success
narrative.  The plan treats historical Test Instance A as adaptively reused development
evidence and prospective Test Instance B as an outcome-sealed field case.  The
latter's pre-freeze corrections already informed the requirements, so it is not
a wholly held-out architecture evaluation.  For Test Instance B, only
producer-side local preflight has completed: package/hash capture, inspection of
packaged source, configuration, and license material, and three embedded-resource
byte checks.  Upstream provenance and licenses, the metric, judge, data roles,
aggregation, environment, locality, and fixed run command/protocol remain
unresolved.  No native ARGO implementation, comparative run, instance model fit,
external outcome, or performance result exists.  Human authority remains
exclusive for real transmission, spending, release, external evaluation, and
final scientific conclusions.  This manuscript therefore presents a design and
prospective evaluation protocol, not implementation or efficacy.
\end{abstract}

\section{Introduction}

A capable language model is only one part of a long-horizon agent.  The agent
also depends on the machinery that stores state, exposes computation, launches
subtasks, recovers from failure, and decides which evidence enters its active
context.  The Recursive Language Models abstraction treats long context as an
external environment that can be inspected and decomposed programmatically
\citep{rlm2025}.  Prime Agent operationalizes this idea in an open-source
harness with a persistent IPython REPL, recursive sessions, daemon continuity,
recovery, direct agent messaging, and durable harness state
\citep{primeagent2026}.  ARGO begins from that inherited substrate.

The scientific problem addressed in this thesis is narrower than general agent
capability and stricter than task completion.  An autonomous empirical agent
must decide what question is worth asking, distinguish competing mechanisms,
freeze a valid comparison before seeing its outcome, preserve the identity of
the code and artifacts that produced the outcome, and limit every conclusion to
the population and protocol actually measured.  It must also expose the points
at which human authority remains decisive.  A run that produces a favorable metric
without these properties may be useful engineering, but it is not sufficient
evidence for a scientific claim.

\subsection{Related work by harness facet}

The retained literature supports an uneven map of the agent harness rather than
a single, complete precedent.  I organize related work by the facet it informs
and treat every unmapped facet as an evidence gap, not as proof that no prior
system exists.

\paragraph{Tools, protocols, and human interaction.}
The Prime Agent paper and pinned v0.9.1 source establish the inherited
code-mediated tool surface and persistent execution substrate
\citep{primeagent2026,primeagentsource2026}.  EurekAgent broadens the design unit
from a tool list to permissions, artifacts, budgets, and human supervision
\citep{eurekagent2026}.  These sources partially cover \emph{tools}, but
\emph{protocols} evidence remains implementation-centered: the pinned Prime
source contains ACP, RPC, MCP, daemon, and REPL seams, while versioned external
protocol and experiment-lifecycle specifications still require retained
primary-source review.  The \emph{agent--user} and \emph{approval-loop} facets
are thinner.  Focused full reads remain outstanding for human--agent
communication, controlled human-in-the-loop workflows, and progressive
autonomy.  Interruption, uncertainty display, abstention, approval identity,
non-bypass behavior, delay, and user burden therefore remain evaluation
questions rather than ARGO properties.

\paragraph{Agent communication, orchestration, and evaluation.}
Prime Agent documents recursive sessions and direct messaging
\citep{primeagent2026}; \emph{Accelerating scientific discovery with
Co-Scientist} provides a prior example of specialized generation, reflection,
ranking, evolution, and meta-review roles \citep{coscientist2025}; and SciAgents
grounds multi-agent hypothesis generation in a literature-derived graph
\citep{sciagents2024}.  These works inform \emph{agent--agent} communication and
\emph{subagent orchestration}, but they do not establish the value of ARGO's
proposed topology.  The AI Scientist's reported code, control, and review
failures motivate evaluator independence \citep{aiscientist2024}.  EviGraph
already provides an operational typed evidence graph, weak-node repair, and
graph-gated manuscript generation \citep{evigraph2026}.  The \emph{evaluators}
facet still lacks an ARGO experiment, and focused full reads of additional
orchestration and long-horizon evaluator candidates remain outstanding.  Those
candidates name gaps only; they do not support claims in this draft.

\paragraph{Context, memory, knowledge, and reusable procedure.}
Prime Agent and the retained full source for Recursive Language Models motivate
programmatic \emph{working context} and context decomposition
\citep{primeagent2026,rlm2025}.  Two byte-checked RLM locators support this
high-level definition, but its reported empirical results are not evidence of
ARGO efficacy.
Continual Harness provides prior work for reset-free durable changes to prompts,
subagents, skills, and memory \citep{continual2026}.  This is partial evidence
for \emph{compression}, \emph{procedural skills}, and \emph{episodic
experience}; it is not a validated account of safe \emph{personalized memory}.
SciAgents and EviGraph inform \emph{semantic knowledge} and provenance-bearing
graph structure \citep{sciagents2024,evigraph2026}.  Full-read gaps remain for
additional context measurement, reusable-skill, episodic-memory, and
personalization/privacy sources that address consent, correction, deletion,
expiry, and cross-user isolation.

\paragraph{Observability, constraints, and decision operators.}
Prime Agent's transcripts and recovery substrate and EviGraph's typed chains
provide partial prior work for \emph{observability}
\citep{primeagent2026,evigraph2026}.  EurekAgent informs the environment side of
\emph{normative constraints} \citep{eurekagent2026}; The AI Scientist and
Co-Scientist inform automated \emph{decision heuristics} and alternative
selection \citep{aiscientist2024,coscientist2025}.  Release-grade full reads are
still missing for versioned provenance/observability standards, governance and
progressive-authority candidates, and primary planning, search, reflection, and
value-of-information sources.  The current ARGO contracts may define candidate
policies, but they cannot independently prove originality or efficacy.

Retained source artifacts now include raw arXiv TeX for seven cited papers and a
PDF/text fallback for Co-Scientist.  Thirteen exact locators for the present
high-level prior-work claims pass byte, excerpt, and semantic review; future
detailed claims require new locators.  No ARGO harness claim has
reproduced-experiment support.  Discovery-only candidates above are used solely
to name missing work.  The design question is whether supported primitives can
be bound to scientific choice, executable protocol, run receipt, contribution
provenance, and manuscript wording closely enough to reduce invalid transitions.

ARGO treats research autonomy as authority to reason within a contract, not as
permission to act without limits.  The agent may formulate questions,
hypotheses, variables, controls, estimands, metrics, falsifiers, and stop rules.
The thesis author retains authority over external transmission, spending,
risky tools, external evaluation releases, and final conclusions.  This separation is a
methodological requirement: the system is evaluated on whether it respects the
boundary, not merely on whether its proposed action would have been convenient.

\subsection{Ancestry, scope, and current status}

ARGO is the descendant produced by migrating the Prime Agent fork.  The
migration itself is product construction.  Consequently, Prime Agent is neither
a plugin placed inside an already complete ARGO system nor an unnamed source of
ideas.  Prime Agent v0.9.1 at the pinned upstream revision is the current
executable substrate \citep{primeagentsource2026}; the ARGO branch presently
contains planning and research contracts, not a native treatment.  Native
construction remains paused pending complete gates, a successor operational
protocol, and explicit human authorization.

Historical Test Instance A is stopped and frozen.  Its panels and external
outcomes were adaptively inspected while the research process was developed, so
it is bounded development evidence and cannot confirm ARGO efficacy.
Prospective Test Instance B is an outcome-sealed field case.  Its pre-freeze
corrections already shaped this protocol and are development evidence; the case
is not a wholly held-out or sealed evaluation of the architecture.  If its
outcome later triggers an ARGO repair, it becomes development evidence for that
repair and cannot confirm the descendant system.

For Test Instance B, only producer-side local preflight has completed:
package/hash capture; inspection of packaged source, configuration, and license
material; and three embedded-resource byte checks.  Upstream provenance and
licenses, the metric, judge, data roles, aggregation, environment, locality, and
fixed run command/protocol remain unresolved.  No characterization run, model
fit, external experiment run, deployable build, external evaluation action, or
performance result exists.

The current planning pointer is a mutable, hash-stamped record that refers to a
separate content-hashed recovery receipt; the planning record itself is neither
content-addressed nor immutable.  The receipt resolves twelve prior references
semantically and preserves one invalid manifest/archive binding.  Prior records
remain historical rationale.  Without retained snapshots of every record
revision, this manuscript does not claim that their bytes stayed unchanged.

\subsection{Thesis objective}

I investigate whether an inherited long-horizon harness can be migrated into an
evidence-grounded autonomous research system without losing behavioral
continuity or laundering provenance.  The thesis has three linked objects of
study:

\begin{enumerate}
  \item a candidate scientific-decision policy intended to freeze rationale,
  protocol, and interpretation before outcome visibility;
  \item requirements for a future evidence architecture that would bind
  research objects, retained records, code, run receipts, artifacts, claims,
  and build receipts; and
  \item a migration and evaluation method that separates inherited,
  re-derived, modified, candidate-original, and external-oracle facts.
\end{enumerate}

These are candidate contributions until source archaeology, native
implementation, ablation, and held-out evaluation support the required wording.
No change of product name, interface, or programming language is treated as
proof of originality.

\section{Research Questions}

The thesis is organized around the following research questions.

\paragraph{RQ1: Decision validity.}
Can immutable, pre-outcome DecisionRecords improve the validity and auditability
of autonomous research choices relative to the same execution substrate using
ordinary session notes?  The relevant failures include post-hoc changes to the
metric or split, absence of a mechanism-distinct alternative, weak
identifiability, and conclusions that exceed the measured scope.

\paragraph{RQ2: Evidence and comparison identity.}
Can three-part comparison identity and a typed evidence graph reduce and detect
unmatched comparisons, stale-claim reuse, cross-instance leakage, and
claim-to-artifact drift?  The study measures both false acceptance of invalid
comparisons or claims and false rejection of valid, properly scoped ones; it
does not assume prevention.

\paragraph{RQ3: Safe self-improvement.}
Does separating research refinement from engine refinement improve held-out
research quality without allowing the engine to train on, rewrite, or optimize
against its own scientific judge?  Can rollback and noninterference tests detect
when durable harness adaptation preserves a harmful shortcut?

\paragraph{RQ4: Migration with honest provenance.}
Can ARGO preserve the observable behavior of the inherited Prime Agent substrate
while adding authority, evidence, and paper-lineage contracts?  Which mechanisms
remain \texttt{INHERITED}, which are genuinely \texttt{MODIFIED}, and which
candidate \texttt{ORIGINAL\_ARGO} mechanisms survive source comparison and
matched evaluation?

\paragraph{RQ5: Prospective process utility.}
Can the candidate process keep data roles, evaluator isolation, comparison
identity, runtime/locality gates, and outcome interpretation frozen on
prospective Test Instance B?  The field outcome is one-shot and externally
reported.  Because pre-freeze corrections already informed the protocol, this
question concerns process execution on the current case, not a wholly held-out
architecture effect.

\section{Architecture Research Methodology}

\subsection{Unit of scientific work}

The proposed unit of ARGO research is not a chat turn or code diff but a
consequential scientific choice represented by a DecisionRecord.  The candidate
record schema binds the active objective and human authority boundary, the
problem and unresolved gap, alternatives, evidence and counterevidence,
identifiability assumptions, population, unit of analysis, data provenance,
split, intervention, matched control, estimand, metric, aggregation,
uncertainty, falsifier, guardrails, stop rule, and reopen condition.  It also
records expected information under positive, null, negative, and failed
execution; resource and license costs; critique; rationale; remaining
uncertainty; code parent; and the three comparison-identity objects defined
below.

The current policy asks for at least two mechanism-distinct alternatives before
a consequential experiment is selected.  That count is a candidate guard
against proposal fixation, not a universal scientific necessity.  Its benefit,
false-block rate, latency, and resource cost require matched ablation.  The
future native design would freeze a DecisionRecord before its outcome becomes
visible and create a descendant record for a changed design; no native mechanism
currently enforces this behavior.

\subsection{Candidate admission and selection policy}

The following ten-step sequence is a candidate ARGO policy.  It is not claimed
to be universally necessary or optimally ordered.  Evaluation must compare it
with information-matched shorter policies and report false blocks, missed
violations, human burden, latency, and compute cost.

\begin{enumerate}
  \item \textbf{Authority:} check whether the proposed action is within
  delegated authority.
  \item \textbf{Legality and ethics:} block illegal, unlicensed,
  privacy-violating, unsafe, or academically misleading work.
  \item \textbf{Identifiability:} ask whether available observations can
  distinguish the claim from its null and major confounders.
  \item \textbf{Evidence:} apply claim-type admissibility rather than treating
  a citation as general verification.
  \item \textbf{Alternative mechanisms:} when the candidate policy applies,
  compare mechanism-distinct designs rather than cosmetic variants.
  \item \textbf{Comparability:} freeze nuisance invariants and the allowed arm
  differences.
  \item \textbf{Information value:} prefer a valid design expected to reduce
  important uncertainty at acceptable cost across possible outcomes.
  \item \textbf{Independent attack:} have a separate critic challenge leakage,
  isolation, assumptions, falsifiers, and interpretation.
  \item \textbf{Preregistration:} freeze the selected design and decision rule
  before outcome evidence is visible.
  \item \textbf{Assimilation:} update only the measured scope; execution failure
  updates engineering state rather than the scientific hypothesis.
\end{enumerate}

Recency, popularity, citation count, agent confidence, majority vote, or a
larger convenient local metric cannot establish efficacy.  A negative
experiment does not close a method family outside its measured population,
learner, and protocol.  Conversely, a policy that rejects everything is not
scientifically valid; false acceptance and false rejection are both outcomes.

\subsection{Three-object comparison identity}

Protocol comparability is not equality of every hash.  A legitimate experiment
must differ in its intervention and will often differ in code or configuration.
ARGO therefore proposes three linked objects for each contrast:

\begin{description}
  \item[Invariant protocol key.] The fields that define the common comparison:
  population, data roles, split, target, estimand, metric, aggregation,
  environment class, and deployment contract.  Model selector and access
  policy, base prompt, permissions, budget, evaluator revision, and the
  order/time block are included when they are nuisance variables rather than the
  intervention.
  \item[Per-arm receipt.] The exact realized identity of each arm, including
  model access, delivered prompt, permissions and tools, code, configuration,
  dependencies and resources, seed or clean-session repeat, order/time block,
  evaluator revision, and produced artifacts.
  \item[Planned-contrast manifest.] The only differences allowed between arms,
  the intervention each difference implements, and the causal label that the
  contrast may support.
\end{description}

A pair is matched when all declared nuisance invariants agree and the actual arm
differences equal the preregistered allowed differences.  An undeclared
difference makes the causal contrast invalid; a declared treatment difference
does not.  In particular, requiring complete arm receipts or code commits to be
identical would reject the intervention being studied.  A future implementation
may hash each canonical object for identity, but it must compare their typed
fields and planned differences rather than collapse comparability into one
full-receipt equality test.

\subsection{Inherited execution substrate}

ARGO inherits its long-running execution substrate from the Prime Agent paper
and the pinned v0.9.1 source revision
\citep{primeagent2026,primeagentsource2026}.  The inherited inventory includes
the daemon supervisor and resident workers; persistent IPython computation and
programmatic context management; recursive RLM sessions and daemon-mediated
messaging; JSONL transcripts, session artifacts, kernel snapshots, goals,
schedules, and recovery; Continual Harness prompts, memories, skills, subagent
specifications, refinement history, and rollback; and TUI, JSON, RPC, and CLI
clients.  These are implementation and ancestry facts, not ARGO results.

Continual Harness supports reset-free online edits to prompts, subagents, skills,
and memory from trajectory windows \citep{continual2026}.  That mechanism
motivates a durable engine-refine lineage.  The risk evidence is separate: a
Prime Agent evaluation records a specification shortcut that was preserved as a
reusable skill despite an anti-cheating heartbeat \citep{primeagent2026}.  That
observed failure motivates least-privilege actions, independent held-out checks,
and auditable rollback; it is not a result of the Continual Harness paper.  The
full RLM source is now retained \citep{rlm2025}.  Its mechanism can inform
context requirements, but its reported empirical results are not evidence that
ARGO works.

\subsection{Typed research and claim-admissibility graph}

SciAgents grounds hypothesis generation in a literature graph
\citep{sciagents2024}.  EviGraph goes further: it already implements a typed
evidence graph, weak-node repair, and graph-gated manuscript generation
\citep{evigraph2026}.  ARGO therefore does not claim the whole graph or paper
lineage as original.  Its planned, re-derived or modified deltas are narrower:
answered-run nodes that cannot be rewritten; three-object arm-contrast identity;
authority and contribution tags bound to claims; retained source bytes; and an
immutable paper-build receipt.

A future native projection would type objectives, authority boundaries,
problems, prior work, hypotheses, designs, DecisionRecords, protocols, runs,
artifacts, results, counterevidence, claims, and builds.  Legal edges would state
which protocol a run executes, which run produced an artifact, which artifact
contains a result field, and which scoped evidence supports a claim.  New
reviewed descendants could repair an interpretation without changing the bytes
or identity of an answered run.  These are requirements; no native graph
currently provides them.

Evidence admissibility is a claim-type lattice rather than one scalar ladder.
Table~\ref{tab:admissibility} summarizes the minimum evidence class for each
claim.  The classes are orthogonal: satisfying one row does not promote evidence
for another row.

\begin{table}[ht]
\centering
\small
\caption{Claim-type evidence admissibility.}
\label{tab:admissibility}
\begin{tabular}{p{0.20\linewidth} p{0.42\linewidth} p{0.27\linewidth}}
\toprule
Claim type & Admissible evidence & Insufficient by itself \\
\midrule
Literature method or reported result & Full paper with retained bytes and a
claim locator; reported scope must remain explicit. & Discovery metadata,
abstract text, citation count, or a summary without source bytes. \\
Software identity or behavior & Pinned primary code/specification plus
independently checked behavior artifacts for the stated revision. & A paper,
README claim, self-attested field, or unpinned repository name. \\
License or provenance & Authoritative license, upstream provenance, and
redistribution/derivative records bound to exact material. & Public access,
software behavior, a model card alone, or silence in documentation. \\
Local efficacy & Exact run artifacts under the matched comparison objects,
with the frozen estimand and uncertainty analysis. & Literature results,
unit-test success, local prose, or an execution receipt without outcomes. \\
External outcome & Identity-bound receipt from the unnamed external oracle. & A
local result, projection, model confidence, or an external value without artifact
identity. \\
\bottomrule
\end{tabular}
\end{table}

A discovery record may motivate a question but cannot support detailed method
or result wording.  A full paper can support what prior authors report but not
local software behavior.  Primary code can establish identity without proving a
license or efficacy claim.  A local run cannot establish an external outcome,
and an external receipt cannot identify which mechanism caused it.  Boolean
verification fields, transcript summaries, and durable harness memories do not
self-attest any row.

\subsection{Experiment lifecycle boundary}

The local ARGO design assigns an unnamed external service the role of
experiment-lifecycle authority and plans a read-only receipt boundary rather
than a competing lifecycle.  This is a statement about the ARGO contract, not an
independently verified claim about that service.  A pinned external source or
specification receipt has not yet been retained, so exact branch immutability,
command/environment binding, artifact identity, and failure semantics remain
open protocol gaps.

If that external evidence is pinned, the planned importer will bind an
experiment node, commit, command, environment, stdout evidence, and artifacts.
Any later migration of lifecycle capability into ARGO must cross an explicit
capability gate and demonstrate identity against the external receipts.  The
boundary separates whether a run was executed reproducibly from whether it was
a scientifically valid run to execute.  An external lifecycle can supply run
facts; the ARGO decision contract supplies authority, identifiability,
comparison, scope, and claim-admission logic.  Neither system's self-description
substitutes for evidence.

\subsection{Research refinement and engine refinement}

ARGO proposes two independent improvement lineages.

\begin{description}
  \item[Research refinement] would update hypotheses, designs, and scoped
  beliefs through admissible evidence, create descendant research states, and
  leave completed-run protocols unchanged.
  \item[Engine refinement] would update prompts, memories, skills, subagent
  specifications, or native behavior for future work.  Promotion would require
  held-out evaluation, rollback, and a judge not used as its training signal.
\end{description}

The separation is intended to prevent a self-improving harness from making a
scientific hypothesis appear stronger by adapting the measurement process that
judges it.  It remains an \texttt{UNCLASSIFIED} candidate mechanism pending
prior-work comparison, implementation, and held-out evaluation; it is not yet a
validated contribution.

\subsection{Paper as a future downstream evidence view}

EviGraph already demonstrates graph-gated manuscript generation
\citep{evigraph2026}.  ARGO's paper service is therefore a planned
re-derivation/modification, not an original claim to graph-to-paper generation.
A future service would apply the following sequence:

\begin{enumerate}
  \item project an append-only research-event store into a reviewed graph
  snapshot;
  \item close that graph into an evidence snapshot containing exact claim,
  source, protocol, run, artifact, authority, and contribution identities;
  \item lower the snapshot into a typed paper intermediate representation;
  \item run grounding, decision, attribution, numeric, and staleness audits;
  \item render TeX and the bibliography; and
  \item invoke an isolated compiler adapter and write an immutable build receipt.
\end{enumerate}

The specifically planned ARGO deltas are retained source bytes, immutable
answered-run identities, exact arm-contrast linkage, authority/contribution
binding, and the final build receipt.  Empirical values in release mode would be
references to fields extracted from immutable run artifacts rather than numbers
copied from prose.  Raw body TeX, arbitrary citations, and unanchored empirical
values would fail the future typed interface.

Reproducible PDF output would additionally require a pinned compiler and
arguments, fonts, locale, and \texttt{SOURCE\_DATE\_EPOCH}, plus normalized or
suppressed PDF metadata.  Those requirements have not been implemented or
validated, so this thesis does not claim deterministic PDF generation.  This
\texttt{paper.tex} is a private working manuscript written before the native
service exists; it is not evidence that the paper pipeline or immutable build
receipt works.

\section{Provenance and Contribution Boundary}

\subsection{Contribution vocabulary}

Every architecture, method, experiment, and result claim uses one of the tags in
Table~\ref{tab:tags}.  The tag is derived from lineage and evidence, not selected
for rhetorical effect.

\begin{table}[ht]
\centering
\small
\caption{Contribution and attribution vocabulary.}
\label{tab:tags}
\begin{tabular}{p{0.20\linewidth} p{0.31\linewidth} p{0.39\linewidth}}
\toprule
Tag & Meaning & Permitted claim form \\
\midrule
\texttt{INHERITED} & Present in the Prime Agent base before the ARGO branch. &
``ARGO inherits or adopts'' the mechanism, with the Prime Agent paper and frozen base revision. \\
\texttt{RE\_DERIVED} & Independently motivated from literature and revalidated under an ARGO protocol. &
``I re-derived and evaluated the mechanism for ARGO,'' citing prior work and the ARGO experiment. \\
\texttt{MODIFIED} & An inherited component whose contract or behavior changed. &
``ARGO modifies'' a named base behavior, followed by the diff and matched evaluation. \\
\texttt{ORIGINAL\_ARGO} & A mechanism with no equivalent in the base or relevant prior systems. &
``I introduce'' only after source comparison, implementation, ablation, and evidence review. \\
\texttt{EXTERNAL\_ORACLE} & An external system supplies a lifecycle or evaluation fact. &
``The external system reports,'' without implying that ARGO generated the fact. \\
\bottomrule
\end{tabular}
\end{table}

Before evidence supports one of these five tags, a proposed component remains
\texttt{UNCLASSIFIED} candidate.  This withholds attribution; it is not a sixth
contribution category.

\subsection{Initial component ledger}

Table~\ref{tab:ledger} records the current, provisional boundary.  Candidate
originality remains open until source archaeology and experiments are complete.

\begin{table}[ht]
\centering
\small
\caption{Initial ARGO component ledger.  Every non-inherited row is planned and provisional.}
\label{tab:ledger}
\begin{tabular}{p{0.36\linewidth} p{0.24\linewidth} p{0.30\linewidth}}
\toprule
Component & Initial tag & Evidence required for final wording \\
\midrule
Persistent computation and programmatic context & \texttt{INHERITED} & Prime Agent paper, pinned source paths, and base revision \\
Daemon, worker, detach/reattach, and recovery & \texttt{INHERITED} & Pinned base architecture and lifecycle tests \\
Recursive subagents and messaging & \texttt{INHERITED} & Prime Agent and RLM records plus base tests \\
Continual prompts, memories, skills, and subagent specifications & \texttt{INHERITED} & Continual Harness and Prime Agent records plus pinned source \\
Research and evidence graph & \texttt{RE\_DERIVED} or \texttt{MODIFIED} candidate & EviGraph comparison; isolate immutable answered runs, arm identity, and authority/contribution binding \\
Three-object comparison identity & \texttt{UNCLASSIFIED} candidate & Prior-art review, native implementation, compatible and corrupted contrast fixtures \\
External lifecycle receipt boundary & Planned \texttt{MODIFIED} integration candidate & Pinned capability contract and receipt-identity tests \\
Research-refine versus engine-refine separation & \texttt{UNCLASSIFIED} candidate & Prior-work review, disjoint held-out protocol, noninterference, and rollback \\
Claim, run, and paper lineage & \texttt{RE\_DERIVED} or \texttt{MODIFIED} candidate & EviGraph comparison; isolate retained source bytes, exact contrast linkage, and immutable build receipt \\
ARGO TUI and CLI product surface & Planned \texttt{MODIFIED} candidate & Native product diff and usability/recovery evaluation \\
\bottomrule
\end{tabular}
\end{table}

EviGraph's operational graph and graph-gated manuscript generation make a claim
of whole-graph or whole-pipeline originality untenable \citep{evigraph2026}.
The candidate ARGO contribution is the isolated delta, not a renamed aggregate.
Each manuscript claim must eventually resolve to reviewed text, a contribution
tag, admissible evidence for its claim type, exact source or run identities,
the three comparison objects when applicable, code revision, scope,
counterevidence, and review status.  An inherited component may be explained as
part of a coherent architecture, but its origin remains visible.  A written
contract is not an implemented or empirical contribution.

\subsection{Authorship and research instrument}

I am the thesis author and retain responsibility for the research questions,
interpretation, and final manuscript.  Language-model sessions, recursive
subagents, and their traces are research instruments.  Their transcripts may
establish which instrument produced a proposal or draft, but they do not replace
authorial responsibility and do not establish scientific evidence.  Human
release authority is recorded separately from model generation and automated
review.

\section{Prospective Evaluation Plan}

\paragraph{Results status.}
\textbf{No native ARGO implementation, comparative system run, instance model
fit, external outcome, or performance result exists.}  The conditions,
endpoints, and analyses below are prospective requirements.  Version numbers,
source identifiers, policy steps, and equations are definitions or provenance,
not results.  The successor operational protocol, thresholds, stopping rules,
and uncertainty choices have not yet been frozen.

\subsection{Conditions and information-matched controls}

The primary comparison will use four explicitly different conditions after a
native treatment exists:

\begin{description}
  \item[SP.] Stock Prime Agent v0.9.1 at the pinned revision, with the inherited
  runtime and no ARGO validator or executable research graph.
  \item[SP-INFO.] The same stock condition receives a deterministic static
  rendering of every task fact available to the treatment plus the same neutral
  research checklist and output request.  It receives no executable graph
  validation, contrast enforcement, closure enforcement, or receipt importer.
  This condition controls for more context, better organization, and checklist
  prompting.
  \item[ARGO-FULL.] The future native treatment with the preregistered ARGO
  mechanisms enabled.  This condition does not exist at the evidence cutoff.
  \item[SP-ORACLE.] A development-only stock condition may receive read-only
  feedback from a private requirements oracle.  It tests whether a proposed
  contract is coherent before porting.  It is excluded from the primary native
  efficacy estimate and is never labeled ARGO.
\end{description}

Within a matched block, conditions use the same base model revision and access
policy, thinking level, allowed evidence bytes, machine class, network policy,
tool schemas, authority card, completion rule, and time/token/compute ceiling,
except for differences named in the planned-contrast manifest.  Evaluators
receive a treatment-neutral rendering of the scientific decision, evidence
links, and action.  Condition/session identifiers and schema-specific surface
forms are removed; emitting an ARGO form earns no credit.

\subsection{Primary endpoint and estimand}

The primary endpoint is valid episode completion, abbreviated VEC.  It equals
one only when the assigned condition completes the task, or declares unresolved
for a rubric-valid reason, within the frozen budget and produces the correct
scientific action: evidence-linked, matched to the frozen comparison, scoped to
the available evidence, and free of critical violations.  Form completion,
blanket rejection, unwarranted abstention, or a fluent rationale attached to the
wrong decision scores zero.

The primary intention-to-treat estimand is the paired task-level risk difference

\begin{equation}
\Delta_{\mathrm{VEC}} =
\Pr(\mathrm{VEC}=1\mid\mathrm{ARGO\mbox{-}FULL})-
\Pr(\mathrm{VEC}=1\mid\mathrm{SP}).
\label{eq:vec}
\end{equation}

All assigned blocks remain in the estimate under the predeclared rule.
Condition-caused timeout, crash, invalid artifact, or unavailable state scores as
an episode failure while retaining its scientific label as an execution
failure.  Critical violations are reported individually, including authority or
legality bypass, post-outcome protocol mutation, unmatched causal attribution,
self-attested identity accepted against the bytes, execution failure presented
as scientific falsification, over-broad closure, or inherited/external facts
misattributed to ARGO.  Wall time, active time, tokens, tool calls, compute,
cost, mandated approvals, and unplanned human correction are reported
separately rather than hidden in an unvalidated composite.

\subsection{Pilot-informed sample and uncertainty design}

The present design fixes no arbitrary task, instance-family, stochastic-repeat,
rater, confidence, or error constants.  A development pilot will estimate paired
discordance, task/template intraclass correlation, clean-session variance,
censoring/failure rates, and, when human judgment is required, rubric
disagreement.  Before confirmatory outcomes are accessible, the successor
preregistration will derive sample and uncertainty choices from the approved
claim scope, smallest practically important effect, available compute/time and
annotation resources, and the thesis author's accepted costs of false positive,
false negative, and false block errors.  If those inputs cannot support the
intended claim, the claim will be narrowed or reported as underpowered.

The task is the randomization and primary analysis unit.  Fixed clean-session
repeats characterize stochastic variation but do not create new independent
tasks; they are aggregated within task before the primary comparison.  The
planned primary test uses paired condition-label swaps within matched blocks.
Uncertainty uses a hierarchy-conditioned bootstrap over instance family,
template, and task while preserving each task's repeat set.  The comparison with
SP-INFO is a prespecified placebo/mechanism contrast: if the SP difference
disappears against SP-INFO, the evidence supports an information/checklist
explanation rather than native enforcement.  The single product contrast in
Equation~\eqref{eq:vec} is primary; a fixed secondary family follows the
preregistered multiplicity policy, and all other ablations are exploratory.

Before outcomes open, the user must freeze an acceptable resource-overhead
boundary and any noninferiority guardrails.  Without that record, the thesis may
report validity and resource effects separately but may not call the cost
acceptable.  A human-adjudication plan is likewise derived from pilot
disagreement and rubric ambiguity.  A one-rater design is labeled as such and
supports no inter-rater reliability claim.

\subsection{Evaluator ownership and leakage control}

Pilot tasks, bounded reusable development/validation tasks, and one-shot
terminal tasks are disjoint and are never pooled into one efficacy estimate.
The terminal evaluator implementation, answer keys, corruption map, and
condition mapping remain outside the agent-writable workspace.  Where feasible,
they are authored or finalized only after the treatment revision is frozen.  A
fixed query and promotion budget prevents iterative adaptation to the terminal
set.

Opening any terminal set converts it into development evidence for every
subsequent repair.  That set cannot confirm a descendant treatment.  Confirmation
then requires a third preselected instance or sealed packet whose outcomes did
not shape the repair.  Local E3 fixtures are development fixtures only; terminal
E3 corruptions and compatible controls are evaluator-owned.  E4 trigger,
refinement-development, held-out validation, and terminal sets are mutually
disjoint.  Evaluator independence means re-deriving judged facts from bytes or
held-out truth; rerunning the producer's verifier is only a self-consistency
check.

\subsection{Evaluation E1: migration fidelity and authority}

E1 will compare the pinned Prime Agent base with a future migrated ARGO slice.
The suite will cover daemon continuity, detach/reattach, persistent-kernel
recovery, recursive child creation, messaging, session artifacts, compaction and
resume, Continual Harness refinement and rollback, and TUI/JSON/RPC/CLI
compatibility.  Each planned slice must name its inherited owner, compatibility
class, source revision, tests, migration evidence, and rollback.

Evidence will be behavioral parity or a preregistered intentional difference in
the native runtime.  A Python requirements oracle can support cross-language
fixtures but cannot prove the native implementation correct.  Authority probes
will use simulated transmission sinks, a fake payment endpoint, a fake external
evaluation endpoint, and deny-by-default capability fixtures.  They will test
approval identity, cancellation, resume, and non-bypass behavior without ever
attempting real transmission, spending, release, destructive action, or external
evaluation.

\subsection{Evaluation E2: scientific decisions and comparison identity}

E2 will score treatment-neutral scientific decisions, not DecisionRecord form
completion.  On matched tasks, the evaluator receives a normalized rendering of
the chosen action, rationale, evidence edges, scope, and comparison objects with
condition and schema markers removed.  Credit depends on the correct decision
under evaluator-held truth.  Completing every field while admitting an invalid
experiment, blocking a valid one, or taking the wrong scientific action scores
zero.  Record completeness and use of alternative mechanisms are secondary
manipulation checks for the candidate policy.

Adversarial fixtures will alter nuisance fields one at a time: data bytes, role
assignment, split, target, estimand, metric, aggregation, model selector/access,
prompt, permissions, budget, seed, code/configuration, environment, evaluator
revision, deployment identity, or order/time block.  The comparison is valid
when nuisance invariants match and actual arm differences equal the
planned-contrast manifest; code or configuration may legitimately differ when
that is the intervention.  Compatible controls test false rejection, while
unregistered differences, self-attested identities contradicted by bytes,
post-outcome records, scope inflation, and execution failure mislabeled as
science test false acceptance.  Both error directions and their resource costs
are reported.

\subsection{Evaluation E3: evidence graph and paper lineage}

E3 development fixtures will contain known-valid and known-invalid event
bundles.  Negative cases include discovery metadata used for a detailed claim,
missing retained bytes or locators, stale claim text, cross-instance evidence,
contrast mismatch, self-attested identity contradicted by bytes, raw TeX or
citation injection, unanchored empirical prose, and inputs changed after an
evidence snapshot.  Positive cases test whether a valid scoped claim can survive
compaction, session deletion, branching, and a future rebuild.  Passing these
local fixtures is development evidence only.

After the treatment freezes, evaluator owners will create or finalize separate
terminal corruptions and compatible controls outside the agent-writable
workspace.  The future evaluator will inspect event root, graph projection,
evidence snapshot, typed paper representation, TeX, bibliography, compiler log,
and build receipt.  It will not assume PDF hashes match until compiler, fonts,
locale, time epoch, and metadata normalization are pinned.  Opening a terminal
bundle consumes its confirmatory role; any repair requires a third preselected
bundle.  Human release remains necessary even if every future audit passes.

\subsection{Evaluation E4: dual refinement}

E4 separates trigger episodes, refinement-development episodes, held-out
validation episodes, and one-shot terminal episodes.  No episode or outcome is
shared across these roles.  Research refinement may update scoped beliefs only
through admissible run evidence.  Engine refinement may change durable harness
state from the trigger/development roles, but it may not read held-out or
terminal outcomes, rewrite scientific results, or modify their evaluators.

A candidate refinement requires a trigger receipt, scoped change, held-out
same-instance probes, cross-instance noninterference probes, rollback evidence,
and an unused promotion query under the fixed budget.  The matched analysis
compares the refined state with a contemporaneous unrefined control.  Trigger
improvement cannot justify promotion.  If an opened terminal outcome motivates
repair, that outcome becomes development evidence and cannot confirm the
repaired descendant; a third preselected instance or sealed packet is required.

\subsection{Process-only field case: Test Instance B}

\subsubsection{Scope and anonymity}

Test Instance B is an outcome-sealed prospective field case, not a wholly
held-out test of the architecture.  Corrections made before the current evidence
cutoff informed ARGO's requirements and are development evidence.  The public
manuscript therefore evaluates only whether the proposed process can freeze and
respect data roles, evaluator isolation, comparison identity, runtime/locality
gates, and outcome interpretation.  Application-specific questions, mechanisms,
data descriptions, labels, metrics, and citations remain in a confidential,
identity-bound operational protocol outside this manuscript.  The agent may not
initiate the external evaluation.  Any outcome enters only as an unnamed,
identity-bound oracle receipt supplied by the thesis author.

\subsubsection{Current preflight boundary}

Only producer-side local preflight has completed: package/hash capture;
inspection of packaged source, configuration, and license material; and three
embedded-resource byte checks.  This does not establish upstream provenance,
license sufficiency, evaluator parity, or performance.  Upstream provenance and
licenses, the metric, judge, data roles, aggregation, environment, locality, and
fixed run command/protocol remain unresolved.  No model fit, experiment run,
deployable build, external evaluation, or result exists.

A mutable, hash-stamped planning record points to a separate content-hashed
recovery receipt.  The receipt resolves twelve prior references semantically and
preserves one invalid manifest/archive binding.  The planning record is not
called content-addressed or immutable.  Prior records remain historical; absent
retained snapshots for each revision, no claim is made that their bytes were
unchanged.

\subsubsection{Process gates before any run}

A successor operational protocol must be frozen before any run.  At minimum it
must bind:

\begin{itemize}
  \item exact data identity, provenance, license, derivative, redistribution,
  and delivery evidence;
  \item mutually exclusive fit, reusable-development, diagnostic, and one-shot
  terminal roles, with every related-data group confined to one role;
  \item the metric implementation, edge cases, judge, aggregation, materiality
  rule, uncertainty procedure, and outcome-dependent decision rule;
  \item the invariant protocol key, each arm receipt, and the planned-contrast
  manifest;
  \item the fixed command, code parent, environment, permissions, resources,
  runtime/memory limits, offline behavior, and item-locality probes; and
  \item evaluator ownership, terminal answer material, condition mapping, fixed
  query budget, and human external-oracle authority.
\end{itemize}

Reusable development/validation material may support debugging and selection
but cannot enter the terminal estimate.  The one-shot terminal evaluator,
answer material, and condition mapping remain outside the agent-writable
workspace and are opened only after the complete procedure and candidate are
frozen.  A failed execution remains an engineering outcome; it does not become
scientific falsification.  A local process result cannot be projected into an
external outcome.

If the one-shot field outcome triggers any ARGO repair, the entire opened case
becomes development evidence for that repair.  It cannot confirm the repaired
descendant.  Confirmation then requires a third preselected test instance or
sealed packet whose content and outcomes did not shape the repair.  This rule
also prevents the pre-freeze corrections from being presented as evidence from
a wholly sealed architecture case.

\section{Threats to Validity and Limitations}

The architecture is presently more complete as requirements than as native
software.  Contract inspection and private-oracle success cannot substitute for
runtime tests or a matched treatment effect.  Test Instance A is adaptively
reused development evidence.  Test Instance B is outcome-sealed only after
pre-freeze corrections had already shaped this protocol, so it cannot establish
a wholly held-out architecture effect or domain generality.  If its outcome
causes repair, confirmation must move to a third preselected instance or sealed
packet.

An independent critic can share the producer's blind spots, particularly when
both use related model families or summaries.  Separate writable state, hidden
evaluator ownership, treatment-neutral rendering, corruption canaries, and a
fixed query budget reduce but do not eliminate this risk.  Human judgment can
also be inconsistent; rater count and adjudication therefore follow pilot
disagreement rather than an invented constant.

Contribution tags remain provisional.  EviGraph already narrows the possible
claim for graph and manuscript lineage, and further source archaeology may
require other tags to be downgraded.  Claim-typed evidence improves traceability
but cannot guarantee that a retained source is correct or that a valid
experiment asks an important question.  This private draft was written directly
from reviewed contracts, not emitted by a native paper service.  No deterministic
PDF, immutable build receipt, native graph, or typed paper pipeline has been
demonstrated.

\section{Conclusion}

This thesis frames ARGO as a provenance-preserving research architecture built
by migrating, not obscuring, the Prime Agent substrate.  Its candidate deltas
concern three-object comparison identity, immutable answered-run lineage,
authority/contribution binding, dual refinement, retained evidence bytes, and a
future build receipt.  These remain requirements rather than implemented or
validated contributions.

The next valid step is not efficacy reporting.  It is to complete the unresolved
evidence, authority, evaluator, provenance, runtime, and locality gates and then
freeze a self-contained successor operational protocol.  The current contracts
are not yet that frozen protocol.  Until a native treatment and prospective
comparative evidence exist, the manuscript makes no claim that ARGO improves
scientific quality or that either test instance demonstrates performance.

\begin{thebibliography}{99}

\bibitem[Karten et al.(2026)]{primeagent2026}
Seth Karten, Alex L. Zhang, Kevin Thomas, Sebastian M{\"u}ller, Elie Bakouch,
Daniel Auras, Mika Senghaas, Fares Obeid, Konstantin Dunas, Johannes Hagemann,
and Sami Jaghouar.
\emph{Prime Agent: A Self-Improving RLM Harness}.
ArXiv preprint arXiv:2608.23552v1, 2026.
\url{https://arxiv.org/abs/2608.23552v1}.

\bibitem[Prime Intellect(2026)]{primeagentsource2026}
Prime Intellect.
\emph{Prime Agent v0.9.1 Source Tree}.
Pinned upstream commit:
\url{https://github.com/PrimeIntellect-ai/prime-agent/tree/81ae3cb34d27d38ee37f9e205a1e73694993b344}.

\bibitem[Zhang et al.(2025)]{rlm2025}
Alex L. Zhang, Tim Kraska, and Omar Khattab.
\emph{Recursive Language Models}.
ArXiv preprint arXiv:2512.24601v3, 2025.
\url{https://arxiv.org/abs/2512.24601v3}.

\bibitem[Karten et al.(2026)]{continual2026}
Seth Karten, Joel Zhang, Tersoo Upaa, Ruirong Feng, Wenzhe Li, Chengshuai Shi,
Chi Jin, and Kiran Vodrahalli.
\emph{Continual Harness: Online Adaptation for Self-Improving Foundation Agents}.
ArXiv preprint arXiv:2605.09998v1, 2026.
\url{https://arxiv.org/abs/2605.09998v1}.

\bibitem[Lu et al.(2024)]{aiscientist2024}
Chris Lu, Cong Lu, Robert Tjarko Lange, Jakob Foerster, Jeff Clune, and David Ha.
\emph{The AI Scientist: Towards Fully Automated Open-Ended Scientific Discovery}.
ArXiv preprint arXiv:2408.06292v3, 2024.
\url{https://arxiv.org/abs/2408.06292v3}.

\bibitem[Gottweis et al.(2025)]{coscientist2025}
Juraj Gottweis et al.
\emph{Accelerating scientific discovery with Co-Scientist}.
ArXiv preprint arXiv:2502.18864v2, 2025.
\url{https://arxiv.org/abs/2502.18864v2}.

\bibitem[Ghafarollahi and Buehler(2024)]{sciagents2024}
Alireza Ghafarollahi and Markus J. Buehler.
\emph{SciAgents: Automating scientific discovery through multi-agent intelligent
graph reasoning}.
ArXiv preprint arXiv:2409.05556v1, 2024.
\url{https://arxiv.org/abs/2409.05556v1}.

\bibitem[Ren et al.(2026)]{evigraph2026}
Zhenjiang Ren, Ruiji Li, Xujing Zhang, Ziliang Pang, Shuo Ren, and Jiajun Zhang.
\emph{EviGraph: Evidence-Guided Autonomous Research Agents}.
ArXiv preprint arXiv:2608.04738v2, 2026.
\url{https://arxiv.org/abs/2608.04738v2}.

\bibitem[Xin et al.(2026)]{eurekagent2026}
Amy Xin, Jiening Siow, Junjie Wang, Zijun Yao, Fanjin Zhang, Jian Song, Lei Hou,
and Juanzi Li.
\emph{EurekAgent: Agent Environment Engineering is All You Need For Autonomous
Scientific Discovery}.
ArXiv preprint arXiv:2606.13662v2, 2026.
\url{https://arxiv.org/abs/2606.13662v2}.

\end{thebibliography}

\end{document}
